\section{Introdução}\label{sec:introducao}

O desenvolvimento de sistemas de propulsão e veículos para voo hipersónico (velocidades superiores a Mach 5) representa um dos maiores desafios contemporâneos da engenharia aeroespacial. Componentes expostos ao fluxo aerodinâmico extremo, tais como bordos de ataque, bordos de ataque de asas, dutos de admissão e sistemas de proteção térmica (TPS), enfrentam aquecimento atrito severo, resultando em temperaturas operacionais que superam facilmente 1500 \textdegree{}C, acompanhadas de gradientes térmicos acentuados e oxidação agressiva.

Superligas tradicionais à base de níquel ou cobalto atingem seus limites operacionais próximos de 1100 \textdegree{}C, enquanto cerâmicas ultra-refratárias (UHTCs), embora resistam a temperaturas elevadas, frequentemente sofrem com baixa tenacidade à fratura e suscetibilidade ao choque térmico. Nesse contexto, o paradigma das Ligas de Alta Entropia (HEAs), proposto em 2004, revolucionou o design de ligas metálicas ao explorar misturas equiatómicas ou quase-equiatómicas de cinco ou mais elementos principais.

Dentro desse domínio, as Ligas de Alta Entropia Refratárias (RHEAs) --- constituídas por elementos de alto ponto de fusão como Nb, Mo, Ta, W, Ti, Zr, V e Hf --- emergem como candidatas-chave para substituir superligas convencionais em ambientes de temperatura ultra-alta. Com a expansão acelerada dessa área de pesquisa, torna-se essencial quantificar e estruturar o conhecimento acumulado. Este trabalho realiza uma análise cienciométrica detalhada e revisão crítica a partir do corpus consolidado de 192 publicações altamente especializadas (período 2004--2026), mapeando a estrutura de autoria, artigos de alto impacto, citações locais e a base teórica estabelecida pelos pioneiros do campo.
